This project investigated how two-phase inlet conditions affect geothermal steam--water separator predictions, building on an inherited computational fluid dynamics model. An imposed stratified inlet placed liquid on the outer-wall side and steam in the remaining area. During baseline mesh assessment, continuing changes in liquid inventory and pressure exposed the need to assess liquid removal before interpreting separation performance. The work therefore examined numerical liquid removal, an explicit brine outlet and wall-film drainage, supported by mass-balance checks, numerical sensitivity tests and droplet tracking. A later comparison on the same mesh found about 37\% more retained liquid with the split inlet than with a mixed inlet at matched 10,000-iteration endpoints. Both still sent nearly all incoming bulk liquid through the steam outlet, so this was a sensitivity result rather than a demonstrated improvement in separation.

Several useful capabilities were established. Low-feed preparation, inlet ramping and Coupled pressure--velocity treatment reduced continuity residuals to approximately 0.002--0.005 at full feed. However, the reported full-feed state still had excess liquid removal equivalent to 20.81\% of liquid feed. Later wall-film calculations passed selected local balance checks with the bulk equations active. A matched direct-drain test reduced the lower-wall store by 70.32\%, but total film by only 1.70\% relative to the undrained case. Four simple outlet-resistance tests reproduced the analytical pressure--flow relation and passed their numerical checks. These verified specific treatments without establishing a stationary separator operating state. Longer film calculations retained storage, while full-geometry and mesh studies remained sensitive to iteration, starting state and numerical controls.

When unfinished particle tracks were counted as separated, reconstructed efficiency estimates agreed with earlier project estimates within 0.019 percentage points. However, 16.78--33.29\% of injected particle mass had no resolved final destination. Increasing the tracking allowance tenfold more than doubled the small escaped-mass estimate. Agreement in the reported efficiency estimate therefore did not demonstrate reliable particle separation.

The project established liquid-removal and wall-film modelling capabilities and identified the checks needed to interpret them. A complete, stationary, mesh-independent and physically validated separator solution was not established. Further work should combine the verified components with consistent mass accounting, confirmed geometry and downstream conditions, and resolved particle fates before using the model for performance optimisation.
